%!TEX root = ../thesis.tex

%jama -- background and related work chapter

\chapter{Background and Related Work}
\label{ch:related}

\section{Compartmental Epidemic Models}
\label{sec:sir}

%jama
The mathematical modelling of infectious diseases traces back to Kermack and McKendrick's foundational SIR (Susceptible-Infected--Removed) model introduced in 1927 \citep{kermack1927}.
In the SIR framework, the population is partitioned into three compartments, and the dynamics are governed by a system of ordinary differential equations parameterised by a transmission rate $\beta$ and a recovery rate $\gamma$:
\begin{align}
  \frac{dS}{dt} &= -\frac{\beta I S}{N}, \\
  \frac{dI}{dt} &= \frac{\beta I S}{N} - \gamma I, \\
  \frac{dR}{dt} &= \gamma I.
\end{align}
The ratio $R_0 = \beta/\gamma$, known as the basic reproduction number, determines whether an epidemic grows ($R_0 > 1$) or decays ($R_0 < 1$).
The mathematical properties of this threshold condition are thoroughly reviewed in \citet{hethcote2000}, and early epidemic growth phases are characterised in detail in \citet{chowell2016}.

Numerous extensions of the SIR model have been developed to capture additional epidemiological phenomena \citep{hethcote2000}.
The SEIR model adds an Exposed ($E$) compartment representing individuals who have been infected but are not yet infectious, with incubation period $\eta^{-1}$.
The SEIR dynamics during COVID-19 are studied in \citet{he2020}, who characterise the sensitivity of epidemic trajectories to $\beta$ and demonstrate the importance of accurate parameter estimation for reliable forecasting.
The SIRV model incorporates vaccination, and the SIRVB model \citep{ariyaratne2025} adds a breakthrough infection pathway allowing recovered or vaccinated individuals to be reinfected with probability $b$.
The macro model in the present project extends this to SIRVHDB, adding Hospitalised ($H$) and Dead ($D$) compartments with city-specific dynamics.

\citet{ariyaratne2025} applied the SIRVB model to COVID-19 data from the World Health Organization across 242 countries, demonstrating that the breakthrough parameter $b$ crucially determines the equilibrium reproduction number $R_t$ in the post-pandemic phase and that SIRVB outperforms both SIR and SIRV in explaining late-pandemic dynamics.
Their work confirms that the simplest model including both vaccination and breakthrough effects is sufficient to capture the main kinetic features of a multi-year pandemic.
Epidemic parameter estimation across countries was further enabled by the Johns Hopkins COVID-19 dashboard, which aggregated and curated global case data throughout the pandemic period \citep{dong2022}.

\section{Agent-Based Models in Epidemiology}
\label{sec:abm}

%jama
Agent-based models (ABMs) represent an alternative and complementary approach to epidemic modelling.
Rather than tracking aggregate compartment populations, ABMs explicitly simulate the state, behaviour, and interactions of each individual (agent) in the population.
Agents are characterised by attributes such as disease state, location, vaccination status, and behavioural rules such as movement patterns and contact rates \citep{zhang2025}.
A comprehensive recent review of ABM approaches, applications, and future directions is provided by \citet{zhang2025}, who survey models ranging from simple lattice-based contact simulators to large-scale digital-twin city representations.

ABMs offer several advantages over compartmental models.
They naturally accommodate population heterogeneity, including age structure, comorbidities, and behavioural variation.
They can incorporate spatial structure and explicit contact networks, capturing the role of superspreaders, geographic clustering, and mobility corridors.
They produce stochastic outputs that reflect the inherent randomness of disease transmission at the individual level, which is especially important for small populations or in the early stages of an outbreak when the law of large numbers does not apply.
The importance of spatial spread in epidemic ABMs was demonstrated by \citet{merler2015} in the context of the 2014 Ebola outbreak in Liberia, where an agent-based model capturing spatiotemporal heterogeneity explained the effectiveness of ring vaccination and isolation policies that purely aggregate models failed to predict.

A landmark open-source ABM implementation for COVID-19 is Covasim \citep{kerr2021}, which represents contacts through layered networks (household, school, workplace, community) and has been calibrated to real outbreak data in multiple countries.
The present project uses a simpler two-city spatial model, but the Covasim experience motivates the use of agent-level randomness and structured mixing as key features.
Contact matrices inferred from survey data \citep{zhang2021} provide empirical support for the mean-field infection probability model used in the non-visual simulation.

The tradeoff is computational cost.
While a compartmental ODE model can be integrated in milliseconds, an ABM with millions of agents may require hours of computation.
For this reason, ABMs are typically used for detailed regional analysis or model validation rather than for rapid optimisation or control design \citep{polcz2025}.
The approach taken in the present project reflects this tradeoff: the MPC controller operates on the computationally efficient macro model, while the ABM serves as a higher-fidelity simulation environment for evaluation.
Comparing non-pharmaceutical interventions across different epidemic scenarios illustrates that the choice of model resolution significantly affects policy conclusions \citep{peak2017}.

\section{Multi-Scale Epidemic Modelling}
\label{sec:multiscale}

%jama
Multi-scale models couple epidemic dynamics at different levels of resolution.
In the context of this project, the macro level represents population dynamics across two cities using compartmental ODEs, while the micro level represents individual agents within each city.
The key interface is the transmission rate $\beta(t)$: the macro model computes a population-level effective transmission rate that is then used to parameterise individual infection probabilities in the ABM.

This coupling is consistent with the mean-field approximation underlying compartmental models: the per-contact infection probability in the ABM, when averaged over a well-mixed population, should recover the macro-level transmission dynamics.
Deviations from the mean field - due to spatial clustering, behavioural heterogeneity, or stochasticity - are precisely what the micro model is designed to reveal \citep{grundel2022}.

\section{Model Predictive Control for Epidemic Management}
\label{sec:mpc}

%jama
Model Predictive Control (MPC) is an optimisation-based control strategy in which a finite-horizon optimal control problem is solved repeatedly online.
At each time step, the controller receives the current state, solves an optimisation problem over a prediction horizon, implements the first portion of the resulting control trajectory, and then repeats.
MPC naturally handles state and input constraints, making it well suited to epidemic management, where constraints such as maximum hospital occupancy and minimum intervention levels are directly relevant \citep{stoustrup2023,grune2017}.

\citet{esterhuizen2024} present a rigorous MPC formulation for the SEIR model without terminal ingredients.
They prove recursive feasibility and asymptotic convergence to disease-free equilibria from any initial state in the admissible set, provided the prediction horizon is sufficiently long.
\citet{kohler2021} develop a robust and optimal predictive control framework for COVID-19, demonstrating that constraint-aware MPC can effectively bound infection and hospitalisation curves even under model uncertainty.
\citet{sauerteig2022} further tailor MPC design to compartmental epidemic models and validate it on realistic parameter sets.

Multi-region MPC formulations are particularly relevant to the present two-city framework.
\citet{carli2020} design MPC policies for a multi-region epidemic scenario and study how intervention coordination across regions reduces total burden compared to independent local policies.
Similarly, \citet{tsay2021} formulate a state-estimation and optimal control problem for COVID-19 across US regions, coupling data-driven $\beta$ estimation with MPC feedback.
Nonlinear MPC formulations with logic constraints for epidemic management are studied in \citet{morato2020}, while stability and optimal control strategies for COVID-19 epidemic models are analysed in \citet{liu2021}.
Nonlinear control of infection spread using a deterministic SEIR model is addressed in \citet{sharifi2021}.

Recent work has extended MPC to more exotic plant models.
\citet{zhong2025} combine physics-informed neural networks (PINNs) with MPC to handle partially unknown SIR dynamics, achieving data-efficient control without full model knowledge.
\citet{kosyanov2023} apply MPC to the SIIR model (with two infected classes distinguishing symptomatic and asymptomatic) and compare performance across different prediction horizons.
\citet{sayedhassen2025} design a data-driven MPC for SARS-CoV-2 in Mauritius, demonstrating the applicability of the framework to island epidemic scenarios with constrained healthcare capacity.

The macro model in the present project uses a similar MPC formulation applied to the extended SIRVHDB model with two cities and a travel matrix, with hospital load $H/H_{\mathrm{cap}}$ as the controlled output.
The MPC generates time series $\beta_1(t)$ and $\beta_2(t)$ representing the effective transmission rates in each city under the prescribed intervention policy \citep{mathworks2025}.
A mean-field perspective on epidemic modelling and control that connects the population-level ODE formulation to the individual-level ABM is provided by \citet{roy2023}.

\section{Hybrid and Data-Driven Approaches}
\label{sec:hybrid}

%jama
A growing strand of the literature combines the computational efficiency of compartmental models with the behavioural richness of agent-based simulation.
\citet{polcz2025} develop a hybrid framework that runs MPC on an ODE model for fast online optimisation, then evaluates candidate policies on an ABM for realistic impact assessment.
A similar philosophy motivates the present thesis: the macro ODE drives the MPC, and the ABM provides the evaluation environment.

Hybrid intelligent models combining epidemic simulation with machine learning for COVID-19 prediction and control are studied in \citet{plank2021}, demonstrating the value of complementary approaches.
The EpidemiOptim toolbox \citep{colas2021} provides a unified framework for optimising control policies in epidemic models using reinforcement learning and evolutionary computation, highlighting the algorithmic diversity available for epidemic intervention design.

Control of socially structured compartmental models under parameter uncertainty is addressed in \citet{albi2021}, who study robust intervention design when contact matrix entries are known only approximately.
This work is directly relevant to multi-city epidemic control, where inter-city contact rates are inherently uncertain.
Data-driven approaches for COVID-19 modelling and forecasting are surveyed in \citet{ramazi2021}, providing context for the parameter calibration challenges faced in the present project.

\section{COVID-19 Kinetic Parameters}
\label{sec:params}

%jama
Calibration of epidemic models to COVID-19 requires estimates of several key parameters.
The basic reproduction number $R_0 \approx 3.0$ corresponds to the original SARS-CoV-2 strain without intervention, consistent with estimates reported in the literature \citep{darienzo2020}.
The recovery rate $\gamma = 1/12\ \mathrm{day}^{-1}$ corresponds to an average infectious period of 12 days.
The hospitalisation rate $h = 0.05$ (5\% of infected) and hospital mortality rate $\mu_h = 0.02$ are consistent with moderate estimates from the pandemic period.
Vaccination provides partial protection characterised by an infection breakthrough factor $\eta = 0.6$.
These parameters, used in both the macro and micro models, ensure that the two model levels operate in a comparable epidemiological regime.
